% !PRINTED 75
\weberpage{75}
%<<CONT

$m$ gerade \qquad \textbf{gerade Funktionen}

\[
\begin{aligned}
\Theta_{00}^{(m)}(v) &= F^{\left(\frac{m}{2}\right)}(\Theta_0,\Theta_1)\\
\Theta_{01}^{(m)}(v) &= \vartheta_{00}(v)\,\vartheta_{01}(v)\,F^{\left(\frac{m}{2}-1\right)}(\Theta_0,\Theta_1)\\
\Theta_{10}^{(m)}(v) &= \vartheta_{00}(v)\,\vartheta_{10}(v)\,F^{\left(\frac{m}{2}-1\right)}(\Theta_0,\Theta_1)\\
\Theta_{11}^{(m)}(v) &= \vartheta_{10}(v)\,\vartheta_{01}(v)\,F^{\left(\frac{m}{2}-1\right)}(\Theta_0,\Theta_1)
\end{aligned}
\]

\weq{16}{
\begin{aligned}
&\text{\textbf{ungerade Funktionen}}\\
\Theta_{00}^{(m)}(v) &= \vartheta_{00}(v)\,\vartheta_{10}(v)\,\vartheta_{01}(v)\,\vartheta_{11}(v)\,F^{\left(\frac{m-4}{2}\right)}(\Theta_0,\Theta_1)\\
\Theta_{01}^{(m)}(v) &= \vartheta_{10}(v)\,\vartheta_{11}(v)\,F^{\left(\frac{m}{2}-1\right)}(\Theta_0,\Theta_1)\\
\Theta_{10}^{(m)}(v) &= \vartheta_{01}(v)\,\vartheta_{11}(v)\,F^{\left(\frac{m}{2}-1\right)}(\Theta_0,\Theta_1)\\
\Theta_{11}^{(m)}(v) &= \vartheta_{00}(v)\,\vartheta_{11}(v)\,F^{\left(\frac{m}{2}-1\right)}(\Theta_0,\Theta_1)
\end{aligned}
}

$m$ ungerade \qquad \textbf{gerade Funktionen}

\[
\begin{aligned}
\Theta_{00}^{(m)}(v) &= \vartheta_{00}(v)\,F^{\left(\frac{m-1}{2}\right)}(\Theta_0,\Theta_1)\\
\Theta_{01}^{(m)}(v) &= \vartheta_{01}(v)\,F^{\left(\frac{m-1}{2}\right)}(\Theta_0,\Theta_1)\\
\Theta_{10}^{(m)}(v) &= \vartheta_{10}(v)\,F^{\left(\frac{m-1}{2}\right)}(\Theta_0,\Theta_1)\\
\Theta_{11}^{(m)}(v) &= \vartheta_{00}(v)\,\vartheta_{01}(v)\,\vartheta_{10}(v)\,F^{\left(\frac{m-3}{2}\right)}(\Theta_0,\Theta_1)
\end{aligned}
\]

\weq{17}{
\begin{aligned}
&\text{\textbf{ungerade Funktionen}}\\
\Theta_{00}^{(m)}(v) &= \vartheta_{01}(v)\,\vartheta_{10}(v)\,\vartheta_{11}(v)\,F^{\left(\frac{m-3}{2}\right)}(\Theta_0,\Theta_1)\\
\Theta_{01}^{(m)}(v) &= \vartheta_{00}(v)\,\vartheta_{10}(v)\,\vartheta_{11}(v)\,F^{\left(\frac{m-3}{2}\right)}(\Theta_0,\Theta_1)\\
\Theta_{10}^{(m)}(v) &= \vartheta_{00}(v)\,\vartheta_{01}(v)\,\vartheta_{11}(v)\,F^{\left(\frac{m-3}{2}\right)}(\Theta_0,\Theta_1)\\
\Theta_{11}^{(m)}(v) &= \vartheta_{11}(v)\,F^{\left(\frac{m-1}{2}\right)}(\Theta_0,\Theta_1).
\end{aligned}
}

Daß in dieser Form \gesperrt{alle} $\Theta$-Funktionen darstellbar sind,
ergibt sich auf Grund von § 20 aus folgenden drei Erwägungen.

1. \gesperrt{Zwischen geraden und ungeraden Funktionen kann keine
lineare Abhängigkeit bestehen, wenn nicht schon zwischen den
geraden Funktionen für sich oder den ungeraden für sich eine
lineare Abhängigkeit besteht.}
